% Include after the Methods section. The preamble needs amsmath for align.
\section{Testing the Dust-Only Hypothesis Using Reddening-Insensitive Indices}
\label{sec:dust-only}

This paper's analysis pipeline (Section~\ref{sec:extinction}) explicitly corrected for color changes resulting from wavelength-dependent extinction from Milky Way dust, but it did not correct for extinction resulting from dust within the two target galaxies. This is because adequately correcting for intragalactic dust requires additional data and several assumptions about the spatial distributions of dust and stellar populations within the galaxy, the optical depth of the dust, and its wavelength-dependent extinction properties \cite{SalimNarayanan2020}. For example, Verley et al.\ (2009) estimated the radial distribution of extinction within M33 using the ratio of total infrared to far-ultraviolet luminosity, obtained from Spitzer and GALEX observations, respectively, whilst making assumptions about the relative attenuation of radiation from stellar populations and ionized gas \cite{Verley2009}.

Nevertheless, this paper is interested in whether the obtained NGC 6503 and M33 color indices are \textit{consistent} with a galactic model where radial color variation arises solely from dust, not from stellar population variation. It does this by defining a reddening-insensitive index $Q$ for each galaxy and computing this index in each annulus for which both constituent color indices are available. This methodology was originally developed by Johnson and Morgan (1953) for stellar photometry and has since been applied to the study of stellar population gradients in spiral galaxies \cite{JohnsonMorgan1953,GonzalezGraham1996}.

For each galaxy, the reddening-insensitive index is defined as
\begin{equation}
    Q \equiv C_{\rm blue} - \alpha C_{\rm red},
\end{equation}
where $C_{\rm blue}$ and $C_{\rm red}$ are the color indices corresponding to the shorter and longer wavelength pairs of filters, respectively, and $\alpha$ is the ratio of the color excesses induced by dust in these two indices under an assumed extinction law. Reddening-insensitive indices of this form have previously been defined for Sloan photometric bands \cite{Abazajian2009}.

In particular, given three filters $a,b,c$ of increasing effective wavelength, we can express the observed color indices as
\begin{align}
    (a-b)(R) &= (a-b)_0(R) + A_a(R) - A_b(R),\\
    (b-c)(R) &= (b-c)_0(R) + A_b(R) - A_c(R),
\end{align}
where the subscript $0$ denotes the intrinsic color index, $R$ denotes the galactocentric radius, and $A_\lambda(R)$ denotes the intragalactic extinction in filter $\lambda$ ($a,b$ or $c$) at radius $R$.

Assuming a foreground-screen dust geometry and a spatially invariant extinction law, we can express
\begin{equation}
    A_\lambda(R) = k_\lambda E(B-V)(R),
\end{equation}
where $E(B-V) \equiv (B-V)-(B-V)_0 = A_B-A_V$ is the color excess in the Johnson $B$ and $V$ bands, and
\begin{equation}
    k_\lambda \equiv \frac{A_\lambda}{E(B-V)}
\end{equation}
denotes the ratio of total extinction in filter $\lambda$ to selective extinction between the $B$ and $V$ bands. The values of $k_\lambda$ are determined by the adopted extinction curve, which is commonly parameterized by the ratio of total to selective extinction, $R_V \equiv A_V/E(B-V)$ \cite{Fitzpatrick1999,Schlafly2011}.

Therefore, by choosing
\begin{equation}
    \alpha = \frac{E(a-b)}{E(b-c)}
    = \frac{k_a-k_b}{k_b-k_c},
\end{equation}
we obtain
\begin{align}
    Q(R) &= (a-b)_0(R) - \alpha(b-c)_0(R)\\
    &\quad + \left[(k_a-k_b)-\alpha(k_b-k_c)\right]E(B-V)(R)\\
    &= (a-b)_0(R) - \alpha(b-c)_0(R).
\end{align}

Thus, under the assumed extinction law and dust geometry, $Q$ is independent of the amount of dust along the line of sight, even if this amount varies substantially with galactocentric radius.

For NGC 6503, the reddening-insensitive index is therefore
\begin{equation}
    Q_{\rm NGC\,6503} = (g'-r') - 1.734(r'-i'),
\end{equation}
where
\begin{equation}
    \alpha_{\rm NGC\,6503}
    = \frac{3.303-2.285}{2.285-1.698}
    \approx 1.734,
\end{equation}
using the SDSS extinction coefficients tabulated by Schlafly and Finkbeiner (2011, Table 6) for the Fitzpatrick (1999) Milky Way extinction law with $R_V=3.1$ \cite{Fitzpatrick1999,Schlafly2011}. These coefficients incorporate a recalibration of the normalization of Galactic reddening maps; this normalization cancels in the ratio used to compute $\alpha$.

For M33, the reddening-insensitive index is similarly
\begin{equation}
    Q_{\rm M33} = (u'-g') - 0.919(g'-r'),
\end{equation}
where
\begin{equation}
    \alpha_{\rm M33}
    = \frac{4.239-3.303}{3.303-2.285}
    \approx 0.919,
\end{equation}
using the extinction coefficients from the same study \cite{Schlafly2011}. Both values of $\alpha$ come from the extinction law of \citet{Fitzpatrick1999} and are set independently of the observed galaxy color profiles.

We can therefore test the hypothesis that the observed radial color variation in each galaxy arises solely from dust by examining the radial dependence of its reddening-insensitive index. Specifically, under this hypothesis, the intrinsic color indices are independent of radius, so we expect
\begin{equation}
    \frac{dQ}{dR} = 0.
\end{equation}

We evaluate this by fitting a constant and a linear model to $Q(R)$ over the radial interval for which \textit{both} color indices are reliably measured. In computing the uncertainties in $Q$, we account for the covariance between the two color indices, as both are computed using a common photometric band. In particular, by standard uncertainty propagation,
\begin{equation}
    \sigma_Q^2 =
    \sigma_{C_{\rm blue}}^2
    + \alpha^2 \sigma_{C_{\rm red}}^2
    - 2\alpha\,{\rm Cov}(C_{\rm blue},C_{\rm red}).
\end{equation}

A statistically significant deviation from a constant value of $Q$ would suggest that radial variation in dust extinction alone, under the assumed extinction law, is insufficient to explain the observed color profiles. Conversely, a value of $Q$ that is statistically consistent with a constant would suggest that the observed color profiles are consistent with, though not necessarily attributable to, radial variation in dust extinction alone.

It should be noted that this analysis assumes extinction due to intragalactic dust can be approximated with a spatially invariant, Milky Way-like extinction law. This is a reasonable starting point because Draine et al.\ (2007), in their analysis of 65 nearby galaxies, found that the dust properties of spiral galaxies generally resemble those of the local interstellar medium of the Milky Way, particularly with respect to their dust-to-gas ratios and polycyclic aromatic hydrocarbon abundances \cite{Draine2007}. More specifically, Wang et al.\ (2022) obtained extinction curves along individual lines of sight within M33 using photometric observations of supergiant stars. They found an average extinction curve similar to that of the Milky Way, with $R_V \approx 3.39$, though individual lines of sight exhibited substantial variation, with $R_V$ ranging from approximately 2 to 6 \cite{Wang2022}. In the absence of an independently determined extinction law for NGC 6503, we adopt the same Milky Way-like extinction law as a reference model for both galaxies.

However, the actual effective attenuation law may differ due to differences in dust grain composition, optical depth, and the relative spatial distributions of stars and dust \cite{SalimNarayanan2020}. This is particularly relevant to this paper because the photometry of each annulus integrates radiation from numerous stellar populations subject to potentially different amounts of extinction. Thus, the reddening-insensitive indices serve primarily as a means of testing a simplified dust-only model, rather than as a means of directly recovering the intrinsic color profiles of the two galaxies.

\subsection{Results of reddening-insensitive index fitting}
\label{sec:dust-only-results}

We retained annuli with a flux detection above $3\sigma$ in each of the three bands. This left 26 annuli spanning $0.09$--$4.76$\,kpc for NGC~6503 and 21 spanning $0.49$--$7.03$\,kpc for M33. For M33 we adopted a fractional $u'$ master-flat uncertainty of $0.506\%$ per detector pixel. The left panels of Figures~\ref{fig:dust-only-ngc6503} and \ref{fig:dust-only-m33} show the measured $Q$ alongside the constant and linear fits. The fits include correlations between annuli and assume a fixed radial pattern for each systematic error. Under this assumption, the constant fits gave $\chi^2=158.7$ for 25 degrees of freedom in NGC~6503 and $\chi^2=40.6$ for 20 in M33. These statistics are sensitive to the assumed correlations because an error that moves two annuli together can contribute little uncertainty to their difference. The fitted line in M33 rises slightly even though the measured points generally fall. Its slope is sensitive to the shared uncertainty in the $u'$ color calibration.

To check how the results depend on the assumed pattern of the systematic errors, we varied the flux in each annulus within its estimated sky bounds, including the master-calibration bounds for NGC~6503, and recomputed the allowed range of $Q$. The right panels show these ranges. A single value of $Q$ can explain all the annuli within these bounds if their allowed ranges overlap.

For NGC~6503, the systematic bounds alone allow a common value of $Q$ between $-0.201$ and $-0.103$\,mag. The data are therefore compatible with the dust-only model within these bounds. In M33 the ranges remain separated by at least $0.13$\,mag even after expanding the flux bounds by $2\sigma$ in each band and varying the shared $u'$ color coefficient within $2\sigma$. This suggests that changing the amount of dust alone in the adopted screen is insufficient to explain the M33 profile. The result is provisional because most of the M33 annuli are bluer than the stars used to calibrate $u'$. A change in the stellar populations or a departure from the assumed dust screen could produce the mismatch.

\begin{figure*}[t]
    \plotone{figures/dust_q_ngc6503.pdf}
    \caption{Reddening-insensitive index $Q_{\rm NGC\,6503}$ in the 26 annuli with detections in all three bands. (a) Measured values with $1\sigma$ error bars and the constant (dashed) and linear (solid) fits using the full covariance. The linear fit has $\chi^2=152.9$ for 24 degrees of freedom. (b) The shaded envelopes show the $Q$ values allowed by the sky and master-calibration bounds at each radius, with a lighter band adding $2\sigma$ of measured flux uncertainty in each filter. The green strip marks the interval shared by all annuli after this expansion.}
    \label{fig:dust-only-ngc6503}
\end{figure*}

\begin{figure*}[t]
    \plotone{figures/dust_q_m33.pdf}
    \caption{Reddening-insensitive index $Q_{\rm M33}$ in the 21 annuli with detections in all three bands. (a) Measured values with $1\sigma$ error bars and the constant (dashed) and linear (solid) fits using the full covariance. The linear fit has $\chi^2=27.6$ for 19 degrees of freedom. (b) The shaded envelopes show the $Q$ values allowed by the sky bounds at the nominal $u'$ calibration, with the lighter band adding $2\sigma$ of measured flux uncertainty in each filter. The horizontal lines mark the highest lower endpoint and the lowest upper endpoint of the expanded intervals. They are separated by $0.35$\,mag at the nominal calibration.}
    \label{fig:dust-only-m33}
\end{figure*}
